% Pursuing Stacks (cotes 134-2 to 134-8): archive cote 134-6, site page 35 (PDF page 36), typescript page 353.
% typescript. Draft compared: ps5.tex:3058-3118.
% AI-transcribed by Claude Opus 5.5 (claude-opus-5-5, adaptive thinking, effort high), Message Batches API, batch ps3 (msgbatch_01XdFmuVVBAcgc7EFxdtCqsj), 2026-11-12; not yet checked by eye.
\scanpage{134-6/35}
\mspage{\hand{353}}
or equivalently,  if A is non empty and any object of A has a section
(over $e_{A^\wedge}$ ).  A topos is called \uline{horizontal} if the family of all subob-%
jects of the final object in the category of sheaves \hand{$\mathcal{A}$} is generating.

For instance,  the topos associated to a topological space is hori-%
zontal -  in particular,  an ``empty'' topos is horizontal.  A topos is
both horizontal and vertical iff it is a (2-)final topos,  i.e. equiva-%
lent to the topos defined by a one-point topological space  (i.e. the
category of sheaves is equivalent to (Sets)).

The property of verticality, I feel,  is of interest in it's own
right,  as exemplified notably by lemma 2 p.268 (which we used yester-%
day),  and the related proposition of section 60 (p.149).  It does not
seem at all subordinated to notions such as \inserted{total} asphericity or total
0-connectedness,  and goes in an entirely different direction  - thus
the terminology  ``\uline{strictly} totally aspheric (or totally 0-connected'') ,
which I still used yesterday (hesitatingly\hand{,} I should say),  is definite-%
ly inadequate.  I would rather say  ``totally aspheric (or totally 0-%
connected\hand{)} \uline{and} vertical '' .

\xout{This reminds me of the}  Another point is about the termihology
of \uline{totally aspheric} and \uline{locally aspheric} objects in a category $A^\wedge$
(with respect to a given basic localizer $\underline{W}$),  introduced in section
79 (p.250),  and still used yesterday.  This terminology does not seem
unadequate by itself,  I introduced it because it struck me as sugges-%
tive (and the notions it refers to do deserve a name,  in order to be
at ease).  The trouble here is that \xout{they}\inserted[typed]{it} conflicts with another possible
meaning,  in accordance with the principle insisted upon forcefully
in the reflections of section 66  -  namely that for objects or arrows
in (Cat) \hand{,} or within a category $A^\wedge$,  the terminology used for naming
properties for these should be \xout{such, that} in accordance with the
terminology used for the corresponding topoi or maps of topoi.  Now
\xout{for} we do have already the notions of \inserted{a} locally aspheric and totally as-%
pheric topos,  which therefore should imply automatically the meaning
of these notions for an object of (Cat)  (which was done satisfactorily
months ago)\hand{,} or for an object of a category $A^\wedge$ .  But in the latter
case, there is definitely conflict with the \xout{use of} terminology introdu-%
ced on p. 252.  This conflict has not manifested itself yet in any
concrete situation,  while the \xout{improper}\inserted[typed]{unorthodox} terminology  has been used
quite satisfactorily a number of times.  T\hand{\unsure{h}}erefore,  I would like to
keep it,  as long as I am not forced otherwise.

x x x

\margin{35}
